\documentclass[11pt]{article}
\pdfminorversion=7
\usepackage[T1]{fontenc}
\usepackage{lmodern}
\usepackage{amsmath,amssymb,booktabs,graphicx,geometry,hyperref,xcolor,array,enumitem,listings}
\usepackage[most]{tcolorbox}
\geometry{margin=25mm}
\hypersetup{colorlinks=true,linkcolor=blue!55!black,urlcolor=blue!55!black,citecolor=blue!55!black}
\newcommand{\src}[1]{\texttt{#1}}
\newcommand{\ma}[1]{\mathbf{#1}}
\newcommand{\ve}[1]{\boldsymbol{#1}}
\newcommand{\E}{\operatorname{E}}
\newcommand{\Cov}{\operatorname{Cov}}
\newcommand{\Var}{\operatorname{Var}}
\setlist{itemsep=2pt,topsep=4pt}
\lstset{language=Fortran,basicstyle=\ttfamily\small,columns=fullflexible,keepspaces=true,
  frame=single,rulecolor=\color{black!25},backgroundcolor=\color{black!3},xleftmargin=4pt,xrightmargin=4pt,
  commentstyle=\color{black!55},keywordstyle=\color{blue!45!black},breaklines=true,showstringspaces=false}
\newtcolorbox{question}[1]{colback=blue!3!white,colframe=blue!35!black,boxrule=0.5pt,arc=2pt,
  left=6pt,right=6pt,top=4pt,bottom=4pt,fonttitle=\bfseries,title=#1}
\newtcolorbox{recommend}{colback=green!4!white,colframe=green!35!black,boxrule=0.5pt,arc=2pt,
  left=6pt,right=6pt,top=4pt,bottom=4pt,fonttitle=\bfseries,title=Our recommendation}

\title{SelAction: open questions on the model\\[4pt]
\large Addendum: two further points on the inputs\\
for Piter Bijma and Jack Dekkers}
\author{Austin Putz}
\date{6 October 2026}

\begin{document}
\maketitle

\begin{abstract}
This short addendum follows the open questions sent on 5 October 2026. While mapping every question the program asks and which of our tests reaches it, we found two further points we cannot settle from the code alone. Question~8 is new: the program does not accept a common-environmental effect of zero in the progeny test. The second point adds to Question~6: discrete and overlapping generations treat half-sib groups differently when each sire has one dam. Neither changes an equation, and we have changed nothing in the program.
\end{abstract}

%======================================================================
\section{Question 8: zero common-environmental effect in the progeny test}\label{sec:q8}

When common environmental effects are switched on (\src{initc = y}) and progeny groups are used, the program asks for a separate $c^2$ for the progeny test of each trait (\src{ccprog}). It rejects zero:

\begin{lstlisting}
read *,ccprog(q)
if (ccprog(q).le.0) then
  print *," wrong input, com.env.effect must be higher than 0!"
\end{lstlisting}

The ordinary $c^2$ of each trait (\src{cc}) is checked with \src{cc(p).lt.0} instead, so zero is accepted there. The same check appears in all four routines: \src{sel1s}, \src{sel2s}, \src{sel3s} (\src{seldiscrete.f90}) and \src{ovlp} (\src{selovlp.f90}). It is also in the original code.

The value only enters as a variance,
\begin{equation*}
\sigma^2_{c,\mathrm{prog}} = c^2_{\mathrm{prog}}\,\sigma^2_P,
\qquad
\Cov_{c,\mathrm{prog}}(p,q) = r_c(p,q)\,\sigma_{c,\mathrm{prog}}(p)\,\sigma_{c,\mathrm{prog}}(q),
\end{equation*}
and is never divided by. When common environmental effects are switched off, the program itself sets $c^2_{\mathrm{prog}}=0$ and uses progeny groups normally. So zero appears to be numerically safe. As it stands, a user who wants common environmental effects among full sibs but none in the progeny test (for example, progeny reared in many separate environments) cannot enter that, and must give a small positive value instead.

\begin{question}{Question 8}
Is there a reason the progeny-test $c^2$ must be strictly positive, or should zero be allowed, as it is for the ordinary $c^2$?
\end{question}

\begin{recommend}
Allow $0 \le c^2_{\mathrm{prog}} < 1$, the same rule as for the ordinary $c^2$, unless there is a modelling reason against it. This is an input check only; no equation changes.
\end{recommend}

%======================================================================
\section{Addition to Question 6: half-sib groups when each sire has one dam}\label{sec:q6}

Question~6 asked whether half-sib sources should be removed at every stage when $M_s = M_d$. Mapping the inputs showed that the two kinds of selection handle this case differently:

\begin{itemize}
\item \textbf{Discrete generations} (\src{sel1s}, \src{sel2s}, \src{sel3s}) always ask ``use half-sib groups?''. If $M_s = M_d$, they then remove half-sib sources from the index with the note ``half-sib information is thus not possible'' (\src{note\_matrat}). As described in Question~6, the removal only applies to the final-stage index.
\item \textbf{Overlapping generations} (\src{ovlp}) ask about half-sib groups only if $M_s < M_d$. With $M_s \ge M_d$, the question is skipped, no half-sib group exists, and there is no note.
\end{itemize}

The overlapping-generation rule is the stricter of the two, and it also covers $M_s > M_d$. That case is not handled separately in discrete generations.

\begin{question}{Addition to Question 6}
Should both kinds of selection follow the same rule? If so, is it ``no half-sib groups unless $M_s < M_d$'' (as in overlapping generations), or ask and then remove with a note (as in discrete generations)?
\end{question}

\begin{recommend}
Use one rule for both, decided together with Question~6. Our planned input checker would then reject half-sib groups when $M_s \ge M_d$, with a clear message, instead of accepting them and dropping them later.
\end{recommend}

%======================================================================
\section*{Source}
\src{fortran/} at commit \src{4e2a245}. The complete list of inputs, with the condition under which each is asked, is in \src{tests/input\_map/README.md}.

\end{document}
